\section{Confluent dilations and Euler differential filters}
\label{stce:section}

The finite-dilation construction has a differential limit that retains
the whole ordinary moment block.  The key point is convergence of the
combined kernels in weighted $L^1$, including arbitrary fixed logarithmic
weights.  Thus the limit is valid at the newly convergent integer moments,
where the individual unfiltered terms need not be integrable.
We use the balanced kernels $C_q$ and scalar Mellin functions $Q_q$ from
Lemma~\ref{stbil:scalar}; these analytic foundations and their shift
derivatives are already in the preserved balanced-Stieltjes report.

\subsection{An exact cube average and weighted confluence}

Let $\mathcal E=x\partial_x$.  For integers $q\ge0$, $r\ge1$ define
\begin{equation}\label{stce:operator}
 \mathcal F_{q,r}(x,a)
 =\prod_{\ell=q}^{q+r-1}(\mathcal E+\ell+1)C_q(x,a)
 =x^{-q}\partial_x^r\bigl[x^{q+r}C_q(x,a)\bigr].
\end{equation}
In particular this is an explicit finite Stieltjes differential expression:
\begin{equation}\label{stce:leibniz}
 \mathcal F_{q,r}
 =\sum_{j=0}^r\binom rj(q+j+1)_{r-j}x^j\partial_x^j C_q.
\end{equation}
Here the finite Stieltjes combination defining $C_q$ is given in
\eqref{stbil:kernel} and its following display.
Only derivatives of generalized Stieltjes functions at the same positive
argument occur.

\begin{theorem}[Confluence in logarithmically weighted Mellin norms]
\label{stce:confluence}
Put $\varepsilon=\log\lambda>0$ and let $H_{q,r}^{\lambda}$ be the
finite-dilation kernel in \eqref{stdil:geometric}.  For every $a>0$,
$M\in\mathbb Z_{\ge0}$ and $0<\sigma<q+r+1$,
\begin{equation}\label{stce:L1}
 \lim_{\varepsilon\downarrow0}\int_0^\infty x^{\sigma-1}
       (1+|\log x|^M)
 \left|\varepsilon^{-r}H_{q,r}^{e^\varepsilon}(x,a)
                  -\mathcal F_{q,r}(x,a)\right|dx=0.
\end{equation}
The same convergence holds locally uniformly in positive shifts $a$.
An exact formula, valid before taking the limit, is
\begin{equation}\label{stce:cube}
 \varepsilon^{-r}H_{q,r}^{e^\varepsilon}(x,a)
 =\int_{[0,1]^r}
       \exp\!\left[-\varepsilon\sum_{j=0}^{r-1}(q+j+1)v_j\right]
       \mathcal F_{q,r}\!\left(xe^{-\varepsilon\sum_jv_j},a\right)
       \,d\mathbf v.
\end{equation}
Moreover $\mathcal F_{q,r}=O(1+|\log x|^{q+1})$ at zero and
$\mathcal F_{q,r}=O(x^{-q-r-1})$ at infinity.
\end{theorem}
\begin{proof}
The defining Stieltjes expression shows that every fixed Euler derivative
has the stated logarithmic bound at zero.  At infinity the normally
uniform Euler--Maclaurin expansion of $C_q$ may be differentiated any
fixed number of times with respect to $x$: take extra terms, extend to a
fixed closed sector in $\Re x>0$, and use Cauchy's formula on discs of
radius proportional to $x$.  In particular an Euler derivative preserves
the inverse-power order of the remainder.  The scalar tail is
\[
 C_q(x,a)=\sum_{k=q}^N
   \frac{(-1)^k\zeta(-k,a)}{k!}
   \left.\partial_z^q(1+z)_k\right|_{z=0}x^{-k-1}
   +O(x^{-N-2}).
\]
Each factor $\mathcal E+\ell+1$ kills the term $x^{-\ell-1}$;
this proves the infinity bound, uniformly on compact positive shift sets.
The second equality in \eqref{stce:operator} follows by the ordinary
Leibniz rule and also proves \eqref{stce:leibniz}.

For a smooth function $f$, the fundamental theorem of calculus gives
\[
 \frac{f(x)-e^{-(\ell+1)\varepsilon}f(e^{-\varepsilon}x)}{\varepsilon}
 =\int_0^1e^{-(\ell+1)\varepsilon v}
   (\mathcal E+\ell+1)f(e^{-\varepsilon v}x)\,dv.
\]
The $r$ factors commute.  Multiplying these identities proves
\eqref{stce:cube}, including every normalization factor.  The integrand
converges to $\mathcal F_{q,r}(x,a)$ uniformly in the cube for fixed $x$.
For $0<\varepsilon\le\varepsilon_0$, its dilated argument lies between
$e^{-r\varepsilon_0}x$ and $x$.  The two endpoint bounds therefore give a
common majorant, of logarithmic type at zero and $O(x^{-q-r-1})$ at
infinity.  Multiplication by $x^{\sigma-1}(1+|\log x|^M)$ makes it
integrable precisely on the asserted open strip.  Dominated convergence
proves \eqref{stce:L1}; the uniform versions of the same estimates prove
the local uniformity in $a$.
\end{proof}

\subsection{Ordinary moments and their first logarithmic derivatives}

Set
\[
 A_{q,r}(u)=\prod_{\ell=q}^{q+r-1}(\ell+1-u).
\]
Taking the confluent limit in the finite-dilation Mellin formula gives
\begin{equation}\label{stce:mellin}
 \int_0^\infty x^{u-1}\mathcal F_{q,r}(x,a)\,dx
 =A_{q,r}(u)\,\pi\csc(\pi u)Q_q(u)\zeta(1-u,a),
 \qquad 0<\Re u<q+r+1.
\end{equation}
Equivalently, integration by parts proves the formula first on
$0<\Re u<q+1$, where all boundary terms vanish, and continuation then
uses the improved decay.  Apparent poles inside the larger strip are
removable in the whole product.

For $k=q+j$, $0\le j\le r-1$, put $d=r-1-j$ and
$S=H_d-H_j$, with $H_0=0$.  Then
\begin{align}
 \int_0^\infty x^k\mathcal F_{q,r}(x,a)\,dx
 &=(-1)^{q+1}j!d!\,
      Q_q(k+1)\frac{B_{k+1}(a)}{k+1}.\label{stce:moment}
\end{align}
\begin{align}
 &\int_0^\infty x^k\log(x/c)\mathcal F_{q,r}(x,a)\,dx\notag\\
 &\quad=(-1)^qj!d!\Bigl[
     \{Q'_q(k+1)-(S+\log c)Q_q(k+1)\}\zeta(-k,a)
       -Q_q(k+1)\zeta'(-k,a)\Bigr],\label{stce:logmoment}
\end{align}
Here $c>0$, $Q'_q$ differentiates its argument, and $\zeta'$ differentiates its
order.  Both integrals are ordinary and absolutely convergent.
For the sign and harmonic coefficient, write $u=k+1+h$ and observe
\[
 A_{q,r}(k+1+h)
 =(-1)^{j+1}j!d!\,h\{1-Sh+O(h^2)\}.
\]
Together with $\pi\csc(\pi(k+1+h))=(-1)^{k+1}/h+O(h)$, this gives
both formulas.  All higher logarithmic identities follow by further
ordinary Mellin differentiation, justified by \eqref{stce:L1} and the
same endpoint bounds.

\subsection{An explicit first-Stieltjes differential identity}

Let $L=\log x$ and evaluate all special functions in the next display at
$x+1$.  Define
\begin{align}\label{stce:gamma1-explicit}
 W(x)={}&3L^2+5L+1-(6L+5)\psi(x+1)\notag\\
 &-(6L+2)x\psi'(x+1)-Lx^2\psi''(x+1)\notag\\
 &-6\gamma_1(x+1)-6x\gamma'_1(x+1)-x^2\gamma''_1(x+1).
\end{align}
Primes on $\gamma_1$ here differentiate its argument.  Direct Leibniz
expansion shows
\[
 W=(\mathcal E+2)(\mathcal E+3)C_1(x,1)
   =x^2\partial_x^2 C_1+6x\partial_x C_1+6C_1.
\]
The moment and logarithmic formulas give the exact triple
\begin{align}
 \int_0^\infty xW(x)\,dx&=\frac1{12},\qquad
 \int_0^\infty x^2W(x)\,dx=0,\notag\\
 \int_0^\infty x^2\log(x/c)W(x)\,dx
       &=-\frac{3\zeta(3)}{8\pi^2}\quad(c>0).
       \label{stce:gamma1-triple}
\end{align}
Indeed $Q_1(2)=1$, $Q_1(3)=3/2$, $\zeta(-2)=0$, and
$\zeta'(-2)=-\zeta(3)/(4\pi^2)$.  The independence of $c$ in the
last identity follows also from the vanishing middle moment.
The tail is $W(x)=-11/(360x^4)+O(x^{-6})$, so weight $x^3$ already
gives logarithmic divergence.  For numerical evaluation one may use
the exact argument-derivative conversions
\[
 \gamma'_1(v)=\zeta'(2,v)+\zeta(2,v),\qquad
 \gamma''_1(v)=-2\zeta'(3,v)-3\zeta(3,v),
\]
which follow by differentiating the Hurwitz Laurent expansion.

There is also an independent signed Laplace representation.  Put
\[
 \nu(t)=1+\frac{2e^{-t}}{1-e^{-t}}
 +\frac{2te^{-t}}{(1-e^{-t})^2}
 +\frac{t^2e^{-t}(1+e^{-t})}{(1-e^{-t})^3}-\frac6t.
\]
Then
\begin{equation}\label{stce:W-Laplace}
 W(x)=\int_0^\infty e^{-xt}\nu(t)\{\log(xt)+\gamma\}\,dt.
\end{equation}
Indeed, with $g(t)=(e^t-1)^{-1}-t^{-1}$ the balanced germ is
$x^z\Gamma(1+z)^{-1}\int_0^\infty t^ze^{-xt}g(t)\,dt$,
normally holomorphic near $z=0$.  Integration by parts sends
$\mathcal E+\ell+1$ to $\ell-t\partial_t$ on $g$, even before
differentiating in $z$.  The two factors therefore give
$2g-2tg'+t^2g''=\nu-1$.  The constant can be discarded because
$\int_0^\infty e^{-xt}\{\log(xt)+\gamma\}\,dt=0$.
Now $\nu(t)=-t^3/360+O(t^5)$ at zero and
$\nu(t)=1-6/t+O(t^2e^{-t})$ at infinity.  Thus Fubini is absolute for
$1<\Re u<4$, and
\[
 \int_0^\infty x^{u-1}W(x)\,dx
 =\Gamma(u)\{\psi(u)+\gamma\}\int_0^\infty t^{-u}\nu(t)\,dt.
\]
This gives a separate quadrature for each identity in
\eqref{stce:gamma1-triple}, without evaluating large cancellations among
Stieltjes functions at large $x$.

\subsection{A positive digamma kernel with elementary and zeta moments}

Define
\begin{equation}\label{stce:positive-K}
 K(x)=2\psi(x+1)+4x\psi'(x+1)+x^2\psi''(x+1)-2\log x-3.
\end{equation}
Then $K=-\mathcal F_{0,2}(x,1)$, and
\begin{equation}\label{stce:K-Laplace}
 K(x)=\int_0^\infty e^{-xt}\mu(t)\,dt,\qquad
 \mu(t)=\frac2t-\frac{t^2e^{-t}(1+e^{-t})}{(1-e^{-t})^3}>0.
\end{equation}
Thus $K$ is strictly completely monotone on $(0,\infty)$.
Its first three ordinary moments and two logarithmic moments are
\begin{align}
 \int_0^\infty K(x)\,dx&=\frac12,\qquad
 \int_0^\infty xK(x)\,dx=\frac1{12},\notag\\
 \int_0^\infty x^2K(x)\,dx&=\frac{\zeta(3)}{2\pi^2}.
       \label{stce:K-moments}
\end{align}
\begin{align}
 \int_0^\infty\log(x/c)K(x)\,dx
     &=-\frac{1+\log(2\pi c)}2,\label{stce:K-log0}\\
 \int_0^\infty x\log(x/c)K(x)\,dx
     &=\zeta'(-1)+\frac{1-\log c}{12},\qquad c>0.
       \label{stce:K-log1}
\end{align}
All are absolutely convergent.  In particular $2K(x)\,dx$ is a
probability measure with mean $1/6$ and second moment $\zeta(3)/\pi^2$.

\begin{proof}
Put $g(t)=(e^t-1)^{-1}-t^{-1}$.  The standard subtracted digamma
integral is $C_0(x,1)=\int_0^\infty e^{-xt}g(t)\,dt$.
Integration by parts sends $\mathcal E+\ell+1$ to
$\ell-t\partial_t$ on the density; all boundary terms vanish.
Consequently $\mathcal F_{0,2}=\int e^{-xt}t^2g''(t)\,dt$,
which proves the displayed formula for $\mu=-t^2g''$.

Its positivity is exactly the classical Lazarevi\'c inequality
$\sinh^3u>u^3\cosh u$ at $u=t/2$.  A short proof makes this input
self-contained.  Since $4\sinh^3u=\sinh3u-3\sinh u$, the coefficients
of $u^3,u^5$ in the difference vanish, while those of $u^{2n+1}$
for $n\ge3$ are positive because
\[
 3^{2n+1}-3>4(2n+1)(2n)(2n-1).
\]
At $n=2$ equality holds.  Increasing $n$ multiplies the right side
by $(2n+3)(2n+2)/((2n)(2n-1))<9$, while the left side becomes
nine times its predecessor plus $24$.  This proves the strict
inequality for every $n\ge3$ and hence for every $u>0$.
Also $\mu(t)=t^3/120+O(t^5)$ at zero and $\mu(t)\sim2/t$ at infinity.
Differentiating its Laplace integral proves strict complete monotonicity.

The first two moments and the logarithmic formulas follow from
\eqref{stce:moment}--\eqref{stce:logmoment}.  The scalar coefficient at
$x^{-3}$ vanishes because $\zeta(-2)=0$, so $K(x)=1/(20x^4)+O(x^{-6})$.
Its Mellin transform therefore extends to $0<\Re u<4$ and is
\[
 -(1-u)(2-u)\pi\csc(\pi u)\zeta(1-u).
\]
At $u=3$ its removable value is $-2\zeta'(-2)=\zeta(3)/(2\pi^2)$.
The nonzero $x^{-4}$ tail makes the next integer moment diverge.
\end{proof}

The hyperbolic inequality used in the positivity argument is classical;
see Yang~\cite{stce:Lazarevic}, equation~(1.1).
No literature-priority claim is made for the Euler operator, its Mellin
action, or that inequality. The additions here are the balanced-kernel
weighted confluence, the resulting explicit ordinary identities, and
their integration with the finite-dilation construction.
